% "{'chapitre':'slci_multiphy','classe':('PSI'),'type':('application'),'titre':'Mise à l\\'eau d\\'un robot sous-marin','source':'Concours Centrale - MP 2019','comp':['B2-02',],'corrige':True}"

%\setchapterimage{fig_00.jpg}
\chapter*{Application \arabic{cptApplication} \\ 
Moissonneuse-batteuse LAVERDA 2350 -- \ifprof Corrigé \else Sujet \fi}
\addcontentsline{toc}{section}{Application \arabic{cptApplication} : Moissonneuse-batteuse LAVERDA 2350 -- \ifprof Corrigé \else Sujet \fi}

\iflivret \stepcounter{cptApplication} \else
\ifprof  \stepcounter{cptApplication} \else \fi
\fi

\setcounter{question}{0}
\marginnote{Concours CCMP -- PSI}





\subsection*{Question de cours}
\question{Donner la strucutre de la chaîne de puissance. On précisera en particulier le nom de chacune des fonctions.}

\question{Donner 3 exemples de composants permettant de réaliser chacune des fonctions.}

\subsection*{Moissonneuse-batteuse LAVERDA 2350}
\marginnote{D'après ressources de Florestan Mathurin.}

\begin{marginfigure}
\includegraphics[width=\linewidth]{fig_01}
\end{marginfigure}

On s’intéresse au système de moissonnage battage
d’une moissonneuse dont on donne une description
structurelle ainsi qu’un extrait de son modèle SysML.
Le céréalier, installé dans la cabine de conduite,
parcourt la surface du champ par bandes parallèles
de la largeur de la coupe de la moissonneuse–
batteuse. La céréale est coupée puis battue, c'est-àdire
que le grain est séparé de son enveloppe (la
balle) et de la paille. Nettoyés des impuretés, les
grains sont transvasés dans la remorque qui roule à
coté de la moissonneuse batteuse. La paille et les
balles sont déposées, en andains, sur les chaumes.

L'actionneur de la prothèse est un moteur à courant
continu alimenté par une batterie rechargeable de 16
Volts. L'énergie mécanique est transmise par un
réducteur de type poulies-courroie suivi d'un système
vis-écrou qui adapte cette énergie mécanique pour la
prothèse (ensemble de liaisons entre le pied artificiel
constitué d'une semelle en fibres de carbone et le
manchon ou tibia artificiel). Des ressorts permettent
d'ajuster également l'énergie mécanique fournie au
pied artificiel. L'effort exercé par les ressorts est
directement relié au couple exercé par l'actionneur.

\begin{center}
\includegraphics[width=\linewidth]{fig_02}
\end{center}

Lors de la récolte, la céréale est coupée, couchée puis rassemblée par la coupe (A) avant d’être acheminée par
le convoyeur (B) vers le batteur (C) pour le battage. Sous l’effet de la pression, une grande partie des grains et
des balles traverse la paroi du batteur et tombe à l’entrée du caisson de nettoyage (E). La paille, une partie des
balles et le reste des grains passent sur les secoueurs (D). Alors que la paille et les balles sont éjectées à
l’arrière de la moissonneuse-batteuse et répandues en andains sur le champ, les grains, plus lourds
descendent dans la paille, traversent les alvéoles puis glissent au fond des secoueurs pour être ramenés à
l’entrée du caisson de nettoyage. Le caisson de nettoyage (E) comprend trois niveaux : le plan supérieur percé
d’alvéoles, l’étage médian (la grille) percé de trous calibrés et un fond plat. Installée au fond du caisson, une
soufflerie génère un vent qui traverse la grille et les alvéoles ; ce vent est destiné à mettre les balles et les
impuretés en suspension et à les expulser hors de la moissonneuse. Les grains après être passés au travers des
alvéoles, passent au travers de la grille puis sont acheminés par une vis d’Archimède dans un réservoir de 8
m3, la trémie (F), en attente de vidange. Lorsqu’elle est pleine, la trémie est vidée dans une remorque en
déployant la goulotte de vidange (G). Cette opération peut être réalisée à l’arrêt ou en pendant le battage.
On donne le diagramme de définition de bloc du modèle SysML correspondant au système de moissonnage.



\begin{center}
\includegraphics[width=\linewidth]{fig_03}
\end{center}

\question{A partir du diagramme de définition de bloc disponible, compléter le diagramme de bloc interne.}

\begin{center}
\includegraphics[width=\linewidth]{fig_04}
\end{center}


